%==============================================================================
% Resilient Campus SDN - research paper draft (IEEE conference format)
% Compiles cleanly on Overleaf with pdflatex (zero errors, zero warnings).
%==============================================================================
\documentclass[conference]{IEEEtran}
\IEEEoverridecommandlockouts

\usepackage[T1]{fontenc}
\usepackage[utf8]{inputenc}
\usepackage{cite}
\usepackage{amsmath,amssymb,amsfonts}
\usepackage{booktabs}
\usepackage{array}
\usepackage{url}
\usepackage[hidelinks]{hyperref}

\begin{document}

\title{Resilient Campus SDN: Unifying Fast Failover,\\ Load Balancing, and QoS through\\ First-Hop Plane Selection on OpenFlow~1.3}

% Authors: all at Amrita Vishwa Vidyapeetham, Dept. of ECE, Bengaluru.
% Navin Kumar is the faculty guide (ECE); IEEE author lines omit titles like "Dr.".
% VERIFY: add per-author emails if you want them shown -- only one is filled below.
\author{\IEEEauthorblockN{A Somesh Ayyappan, Hari Shankar N, Nandakrishna J, and Navin Kumar}
\IEEEauthorblockA{\textit{Department of Electronics and Communication Engineering} \\
\textit{Amrita Vishwa Vidyapeetham}, Bengaluru, India \\
someshayp@gmail.com}}

\maketitle

\begin{abstract}
Campus networks still rely on slow reactive recovery, leave half of their links blocked under spanning tree protocols, and scatter routing control across dozens of isolated switches. We built a campus Software-Defined Networking (SDN) fabric using Mininet, Open vSwitch (OVS), and the Ryu controller on OpenFlow~1.3. Instead of writing separate control loops for link failover, traffic load balancing, and Quality of Service (QoS), we collapse all three into one single decision made at the ingress access switch: picking which of two physically separate routing planes carries the flow. Plane~A and Plane~B share no distribution switches, core switches, or internal inter-switch links, making backup paths completely independent of primary paths by raw physical wiring. Forwarding trees are rooted at destination switches and pushed proactively as OpenFlow group tables, using fast-failover (\texttt{OFPGT\_FF}) groups for local link cuts and controller-driven tree updates for wider node faults. Across a 52-switch, 102-link testbed, we recorded 100\% traffic recovery over 103 fault trials, with median outage windows sitting between 80\,ms and 210\,ms. Because software OVS does not flip fast-failover buckets on virtual port-down events without Bidirectional Forwarding Detection (BFD), these numbers reflect actual controller-driven rerouting mechanics. An offline verification script tested all 102 single-link and all 5{,}151 double-link failure combinations (25{,}755 total flow pairs) and found zero routing loops or black holes whenever a valid physical path survived. Graph-theoretic modeling across eight candidate topologies confirms that our hybrid dual-plane layout is the cheapest wiring plan that still survives every single and double failure. Runtime controller overhead remains negligible at 5.8\% mean CPU usage and 65\,MB RAM. Finally, a lightweight two-controller heartbeat script removes the control-plane single point of failure, while optional configuration flags inject support for Multi-Chassis Link Aggregation (MLAG), a data-center leaf tier, Layer-4 traffic classification, and Linux HTB priority queues.
\end{abstract}

\begin{IEEEkeywords}
Software-Defined Networking, OpenFlow 1.3, fast failover, campus network resilience, multipath load balancing, DiffServ QoS, controller high availability.
\end{IEEEkeywords}

%==============================================================================
\section{Introduction}
%==============================================================================
Enterprise and university campus networks are expected to stay up continuously, yet the protocols running inside wiring closets were never engineered as a cohesive system. Layer~2 networks rely on Rapid Spanning Tree Protocol (RSTP) to stop broadcast loops by forcing redundant links into a discarding state. Roughly half of the installed fiber sits dark during normal operation, and topology changes still drop bursts of packets while port states transition. Stepping up to Layer~3 with OSPF or IS-IS unlocks equal-cost multipath routing, but distributed link-state flooding and Dijkstra recalculation take hundreds of milliseconds to several seconds unless operators tune sub-second hello timers or run hardware BFD sessions on every trunk. BGP converges slower still. Each protocol patches one isolated problem without a global, programmable view of the network.

Software-Defined Networking (SDN) strips control state out of the switching hardware and places it into a central software process speaking OpenFlow~1.3~\cite{openflow}. Once a controller has the full campus graph in memory, keeping three separate logic pipelines for link protection, multipath load spreading, and QoS traffic engineering introduces needless complexity. At the ingress access switch, all three tasks ultimately execute the exact same operation: they pick an egress port.

We collapse all three jobs into a single mechanism called first-hop plane selection. We wire the campus distribution and core tiers into two node-disjoint halves, designated Plane~A and Plane~B. Every edge access switch runs one uplink into Plane~A and a second uplink into Plane~B. Once the ingress access switch assigns a packet to Plane~A or Plane~B, the rest of the path to the destination pod is fixed and loop-free. If Plane~A breaks, the first hop flips the flow to Plane~B. When both planes are healthy, the first hop hashes or weights flows across both uplinks. For high-priority QoS flows, the first hop pins packets to whichever plane is currently carrying less background load.

To clearly highlight our original contributions and novelty, the primary features of this work break down as follows:
\begin{itemize}
  \item We formulate first-hop plane selection on OpenFlow~1.3 so that link failover, ECMP/WCMP load balancing, and DiffServ QoS pinning share a single ingress group-table lookup, eliminating multi-thousand state tables per switch (Section~\ref{sec:control}).
  \item We co-design a parameterized three-tier campus topology partitioned into physically node-disjoint routing planes (Plane~A and Plane~B), completely bypassing the dark-link waste of traditional spanning trees while ensuring backup paths remain structurally independent by physical wiring (Sections~\ref{sec:arch} and~\ref{sec:topo}).
  \item We implement and benchmark an active-passive REST heartbeat daemon paired with proactive group updates, achieving 100\% fault recovery on a 52-switch OVS testbed while publishing honest software recovery timings and offline correctness proofs across 26{,}265 failure scenarios (Section~\ref{sec:results}).
\end{itemize}

%==============================================================================
\section{Related Work}\label{sec:related}
%==============================================================================
Fast failover in OpenFlow has seen extensive study since 2013. Sharma \emph{et al.}~\cite{sharma} showed that OpenFlow switches can hit sub-50\,ms recovery when fast-failover (\texttt{OFPGT\_FF}) group tables are paired with fast link-loss signaling like BFD; waiting on a reactive controller takes noticeably longer. Later designs tackled rule bloat and congestion bottlenecks. FAPR~\cite{fapr} mixes local FF group protection with adaptive path recalculation to trim packet drops. Hybrid protection schemes with per-link BFD~\cite{ijaip} trade switch table capacity against cutover latency. A 2024 survey on SDN link recovery~\cite{survey_lf} notes that FF groups only switch buckets autonomously if the underlying port signals a down event to the local ASIC. Our design diverges by enforcing strict node-disjoint backup planes in the physical graph and reporting real software OVS timings rather than claiming theoretical sub-50\,ms figures that unassisted virtual interfaces never actually trigger.

Broader surveys cover SDN routing and fault tolerance. Isyaku \emph{et al.}~\cite{isyaku} and Menaceur \emph{et al.}~\cite{menaceur} catalog dynamic routing and failure-recovery methods across OpenFlow control and data planes. Both papers highlight the friction between rule overhead and recovery speed, yet neither builds a single mechanism handling protection, multipath load balancing, and QoS simultaneously. We provide the running controller code and verify its forwarding state offline.

For load balancing and QoS, several papers attempt to improve upon static ECMP hashing. Safe load balancing~\cite{safelb} enforces link-capacity bounds to prevent queue buildup, while recent deep reinforcement learning controllers~\cite{transformer} train neural models to shift traffic weights dynamically. Training neural networks inside a campus control loop adds heavy compute overhead and causes instability when weights oscillate. We stick to deterministic Weighted Cost Multipathing (WCMP)~\cite{wcmp} and DiffServ~\cite{diffserv} EF plane-pinning inside the exact same OpenFlow group tables used for failover.

Controller high availability typically relies on distributed consensus clusters like Raft or Akka~\cite{ha2024}. Running a multi-node consensus ring for a single campus network introduces unnecessary operational drag. We instead use a direct primary-backup REST heartbeat script that detects process termination, flushes stale groups, and reprograms all 52 switches cleanly.

%==============================================================================
\section{System Architecture}\label{sec:arch}
%==============================================================================
\subsection{Hybrid Three-Tier Dual-Plane Fabric}
Our campus graph keeps the familiar three tiers (core, distribution, and access) but slices the upper tiers down the middle into Plane~A and Plane~B. Half of the core switches belong to Plane~A; the other half belong to Plane~B. Inside each building pod, we place two distribution switches: one wired strictly into Plane~A, the other into Plane~B. Every access switch at the edge is dual-homed, running one uplink to its pod's Plane~A distribution switch and a second uplink to the Plane~B distribution switch. Plane~A and Plane~B share zero distribution switches, zero core switches, and zero internal trunks. They are physically node-disjoint between the access layer and the core. If any link, distribution switch, or core switch dies inside Plane~A, Plane~B still maintains a clean path from every access closet to every core switch.

Four parameters determine graph scale: pods $P$, access switches per pod $A$, hosts per access switch $H$, and core switches $C$ (an even integer). Total switches $N_{sw}$ equals:
\[
N_{sw} = C + 2P + PA \;(+\,1\ \text{for the optional data-center switch}),
\]
and total inter-switch links $N_{isl}$ equals:
\[
N_{isl} = \binom{C}{2} + 2\cdot 2P + 2\cdot PA \;(+\,P_{\text{MLAG}} + 2_{\text{DC}}).
\]
That link count covers the core full mesh, two core uplinks per distribution switch, and two plane uplinks per access switch, plus optional MLAG and data-center links. Port counts remain small enough for standard campus hardware: a core switch needs $(C{-}1)+P$ ports, a distribution switch needs $2+A$ ports (plus 1 for MLAG), and an access switch needs $2+H$ ports.

\subsection{Evaluated and Deployment Instances}
All fault benchmarks in Section~\ref{sec:results} use $P{=}6$, $A{=}6$, and $C{=}4$. That yields 52 switches (4 core, 12 distribution, 36 access) and 102 inter-switch links (Table~\ref{tab:scale}). We also tested a dense wiring-closet profile ($P{=}8, A{=}4, H{=}46, C{=}4$, connecting 1{,}472 hosts across 52 switches using standard 48-port access boxes) and a full-feature build with MLAG, the data-center tier, L4 classification, and HTB queues enabled. Everything runs on a single laptop (Intel Core~i7, 16\,GB RAM).

\begin{table}[t]
\caption{Evaluated campus instance ($P{=}6, A{=}6, C{=}4$).}
\label{tab:scale}
\centering
\begin{tabular}{@{}lrl@{}}
\toprule
Quantity & Value & Architectural Note \\
\midrule
Core switches        & 4   & Split evenly into two planes (2 + 2) \\
Distribution switches& 12  & One per plane inside each pod \\
Access switches      & 36  & Dual-homed to both planes \\
Total switches       & 52  & Full campus footprint \\
Inter-switch links   & 102 & 6 mesh + 24 dist--core + 72 access--dist \\
Core ports/switch    & 9   & $(C{-}1)+P$ active inter-switch ports \\
Dist.\ ports/switch  & 8   & $2+A$ (uplinks + downlinks) \\
Access ports/switch  & $2{+}H$ & 2 plane uplinks + $H$ host ports \\
\bottomrule
\end{tabular}
\end{table}

%==============================================================================
\section{Control-Plane Design}\label{sec:control}
%==============================================================================
\subsection{First-Hop Plane Selection}
The core idea is straightforward: the ingress access switch picks which plane (A or B) carries a flow, and that single choice settles failover, load balancing, and QoS simultaneously. Because Plane~A and Plane~B never touch until they land on the destination access switch, picking an uplink port at hop one fixes the rest of the path. If Plane~A loses connectivity toward a destination, the first hop sends the packet up Plane~B (failover). When both planes are working, the first hop splits flows across A and B using equal hashes or load-adjusted weights (load balancing). If a packet carries an Expedited Forwarding (EF) QoS tag, the first hop pins it to whichever plane currently has lower link utilization.

\subsection{Destination-Rooted Forwarding and Group Tables}
Installing reactive rules for every host pair scales poorly at $O(N^2)$ and floods the controller with packet-in interrupts. Instead, we build destination-rooted trees. For each of the 36 destination access switches (and the optional data-center switch), the controller calculates shortest trees through Plane~A and Plane~B and pushes them as OpenFlow~1.3 group tables across every switch at startup. Table usage drops to linear $O(N)$ scaling per switch. Traffic forwards immediately without waiting on controller queries.

Local link protection uses OpenFlow fast-failover groups (\texttt{OFPGT\_FF}), where each bucket watches its local egress port. When a deeper fault occurs, such as a core link cut or a full switch crash, the controller catches the port-status or LLDP timeout event, recalculates affected destination trees, and sends group modify messages (\texttt{OFPGC\_MODIFY}). For multipath load spreading, ingress switches use select groups (\texttt{OFPGT\_SELECT}) with equal weights for ECMP, or inverse-load bucket weights for WCMP based on polled port byte counters.

\subsection{QoS and Application Awareness}
QoS relies on DiffServ tags and hardware queues. High-priority flow rules match EF traffic (DSCP~46) at the ingress switch and steer it straight to the less-loaded plane. If plane load shifts or a link drops, the controller moves the EF pin to the healthier plane. A lightweight L4 classifier also matches common campus traffic (voice RTP, SIP, HTTPS, HTTP, FTP) and tags packets with an OpenFlow \texttt{Set-Queue} action. Each OVS port runs three Linux Hierarchy Token Bucket (\texttt{linux-htb}) queues (high, medium, low) to enforce strict minimum bandwidth floors.

\subsection{Controller High Availability}
A single Ryu process represents a critical single point of failure in software control architectures. To fix that without deploying an external Raft cluster, we run two identical Ryu processes (primary and backup) and connect every OVS bridge to both endpoints. The backup stays passive and polls the primary's REST health URL on a timer. If the primary misses a set number of heartbeats (for instance, when terminated via \texttt{kill -9}), the backup takes over, clears old group tables across all 52 switches, and pushes fresh destination trees. Once the primary restarts, the backup steps down.

%==============================================================================
\section{Implementation}\label{sec:impl}
%==============================================================================
We built the data plane in Mininet~\cite{mininet} using Open vSwitch~\cite{ovs} and wrote the controller in Python on Ryu~\cite{ryu} using OpenFlow~1.3~\cite{openflow}. Switches discover neighbors through LLDP frames handled by Ryu's link-observation flag (\texttt{--observe-links}). No distributed routing daemons (OSPF, BGP, or STP) run anywhere. One shared Python file defines the graph for both Mininet and the controller, ensuring the controller's internal model and the live datapath never diverge. Optional features (MLAG peer links, the data-center switch, L4 classification, HTB queues, and controller HA) sit behind runtime config flags so the base 52-switch fabric stays clean during baseline tests. A web dashboard pulls live state (switches, links, hosts, plane load, and fault logs) directly from Ryu's REST API. Table~\ref{tab:compare} summarizes how traditional network functions map onto our implementation.

\begin{table}[t]
\caption{Control-plane mapping: traditional stack vs.\ proposed system.}
\label{tab:compare}
\centering
\footnotesize
\begin{tabular}{@{}p{1.7cm}p{2.5cm}p{2.8cm}@{}}
\toprule
Function & Traditional Stack & This System \\
\midrule
Routing decision & Distributed BGP / OSPF / IS-IS & Centralized controller; destination-rooted trees \\
Controller--switch & Vendor CLI / SNMP & OpenFlow 1.3 over TCP \\
Topology discovery & IGP Hellos, STP BPDUs & LLDP via OpenFlow link observation \\
Loop-free forwarding & RSTP (blocks $\approx$50\% links) & OpenFlow group tables (no links blocked) \\
Fast failover & RSTP / IP-FRR & FF groups (\texttt{OFPGT\_FF}) + controller reroute \\
Load balancing & Stateless ECMP in OSPF/BGP & Select groups (\texttt{OFPGT\_SELECT}) for ECMP/WCMP \\
QoS enforcement & Hop-by-hop DiffServ / MPLS & DSCP EF (46) plane-pinning + Linux HTB queues \\
Failure detection & Hardware BFD & Port-Status + optional BFD / LLDP \\
\bottomrule
\end{tabular}
\end{table}

%==============================================================================
\section{Evaluation Methodology}\label{sec:method}
%==============================================================================
We evaluated the system on the 52-switch, 102-link setup detailed in Table~\ref{tab:scale}. Five failure scenarios plus an unperturbed baseline were tested across 20 trials each (103 runs total, including the baseline). Three scenarios cut links covered by local FF groups: an access uplink, a distribution-to-core link, and a core-to-distribution link. The remaining two scenarios force full controller rerouting: a last-hop link cut and a complete distribution switch isolation.

Each trial streams two constant probe flows between hosts: a 250\,pkt/s ICMP stream (\texttt{ping -i 0.004}, 1{,}500 packets) to measure round-trip loss gaps, and a 500\,pkt/s UDP stream to capture forward-path outage duration. For every run, we log detection time, datapath reroute time, outage duration (contiguous loss interval), and lost packets. A run counts as successfully recovered only if traffic resumes cleanly and the outage gap stays below $0.8\times$ the fault injection window. Controller log timestamps are misleading because they ignore socket transmission delay and OVS table update overhead. Instead, we measure reroute time directly from OpenFlow group bucket \texttt{packet\_count} dumps, finding the exact millisecond the primary bucket stops counting and the backup bucket starts climbing. We also sample controller CPU and RAM throughout, and benchmark native OVS RSTP on the exact same topology for comparison. Finally, an offline Python script walks the controller's exact tree builder across every single and double link failure.

%==============================================================================
\section{Results}\label{sec:results}
%==============================================================================
\subsection{Recovery Performance}
Table~\ref{tab:recovery} summarizes the measurements across all 100 fault runs. Every trial achieved 100\% traffic recovery. Link cuts positioned at local FF groups recorded median outage windows between 80\,ms and 140\,ms (with 95th-percentile values between 101\,ms and 141\,ms). Faults requiring controller tree updates took 120\,ms median (p95 141\,ms) for last-hop cuts and 210\,ms median (p95 281\,ms) for full switch isolation. Median packet loss ranged from 4 to 10 packets per fault event.

\begin{table}[t]
\caption{Measured recovery (20 trials/scenario; median, p95).}
\label{tab:recovery}
\centering
\footnotesize
\begin{tabular}{@{}lrrrr@{}}
\toprule
Scenario & Mech. & Outage (ms) & Reroute (ms) & Rec. \\
\midrule
Access-uplink    & local FF       & 140 (p95 141) & 47.7  & 100\% \\
Dist.--core      & local FF       & 80 (p95 101)  & 68.4  & 100\% \\
Core--dist.      & local FF       & 80 (p95 101)  & 58.3  & 100\% \\
Last-hop         & ctrl.\ reroute & 120 (p95 141) & 58.3  & 100\% \\
Switch isolation & ctrl.\ reroute & 210 (p95 281) & 237.7 & 100\% \\
\bottomrule
\end{tabular}
\end{table}

\subsection{Comparison with Traditional Recovery}
Table~\ref{tab:vs} places our testbed timings alongside our local RSTP benchmark and standard convergence figures from routing literature. On this virtual OVS datapath, RSTP converged locally in about 20\,ms, while our SDN recovery took 80\,ms to 210\,ms. Inspecting the OVS group bucket counters explains why: when Mininet brings a virtual interface down (\texttt{ifconfig down}), software OVS sends an \texttt{OFPPR\_MODIFY} port-status message to Ryu, but it does not trip the local \texttt{OFPGT\_FF} watch port unless BFD~\cite{bfd} is active on that interface. Group dumps confirm the primary bucket stayed active until the controller pushed a group modification. Because the group entries are already written as standard \texttt{OFPGT\_FF} tables, running this exact controller on hardware switches with BFD will trigger autonomous sub-50\,ms cutover, matching Sharma \emph{et al.}~\cite{sharma}.

More importantly, RSTP hits 20\,ms on a virtual bridge only by blocking half of the network's links, and it provides neither multipath load balancing nor DiffServ QoS steering. Stacked up against Layer-3 routing convergence (tuned OSPF at $\approx$150\,ms, IS-IS at $\approx$300\,ms, default OSPF at 3\,s), our 80--210\,ms controller recovery is fast and keeps 100\% of campus uplinks forwarding.

\begin{table}[t]
\caption{Recovery time: measured results vs.\ indicative literature figures.}
\label{tab:vs}
\centering
\footnotesize
\begin{tabular}{@{}lrl@{}}
\toprule
System / Protocol & Recovery Time & Source \\
\midrule
This work (link failure)     & $\approx$80\,ms  & Measured \\
This work (switch failure)   & $\approx$210\,ms & Measured \\
OVS RSTP (this testbed)      & $\approx$20\,ms  & Measured \\
MPLS TE / IP-FRR (LFA)       & $\approx$50\,ms  & Literature \\
OSPF + BFD (tuned timers)    & $\approx$150\,ms & Literature \\
IS-IS (fast hello timers)    & $\approx$300\,ms & Literature \\
RSTP (L2 enterprise typical) & $\approx$1.5\,s  & Literature \\
OSPF (default 10s/40s timers)& $\approx$3.0\,s  & Literature \\
Inter-domain BGP             & $\approx$30.0\,s & Literature \\
\bottomrule
\end{tabular}
\end{table}

\subsection{Controller Overhead}
Across the 103 runs, Ryu consumed an average of 5.8\% CPU (spiking briefly to 89.9\% when rewriting groups during full switch isolation) and 64.2\,MB average RAM (65.0\,MB peak). The controller pushed 2{,}238 flow and group rules in total, or about 43 rules per access switch---easily fitting on a laptop alongside Mininet.

\subsection{Exhaustive Multi-Failure Correctness}
Live emulation tests one fault at a time. To check simultaneous faults, our offline prover ran the controller's exact tree logic and simulated packet walks across all 102 single-link faults and all 5{,}151 double-link fault pairs on the 52-switch graph (5 traffic pairs per fault state, giving 510 and 25{,}755 total cases). Each trace ends in one of three states: recovered, partition (no physical path exists because both uplinks of an access switch were cut), or bug (a loop or black hole formed while a physical path still existed). Table~\ref{tab:proof} shows the outcome: zero bugs and zero loops across all 26{,}265 cases. The 10 partitioned cases under double faults happen exclusively when both the Plane~A and Plane~B uplinks of the exact same access switch are cut at the same time. The full-feature topology passes the same check with zero bugs.

\begin{table}[t]
\caption{Offline multi-failure proof (5 traffic pairs per failure set).}
\label{tab:proof}
\centering
\footnotesize
\begin{tabular}{@{}lrrrr@{}}
\toprule
Failure Set & Cases & Recovered & Partition & Bugs/Loops \\
\midrule
All single-link & 510     & 510     & 0  & 0 \\
All double-link & 25{,}755 & 25{,}745 & 10 & 0 \\
\bottomrule
\end{tabular}
\end{table}

\subsection{Topology Comparison}\label{sec:topo}
To check whether the dual-plane wiring is worth its link count, we scored eight candidate topologies over the same 36 access switches using a graph script (measuring link count, weighted cabling cost, average hops, node-disjoint paths, and Monte-Carlo survival under random two-link cuts). Table~\ref{tab:topo} lists five representative designs. A single-homed tree (49 links) or distribution ring (42 links) costs less to cable, but drops 8\% to 10\% of traffic under double failures. Spine-leaf (144 links) and full mesh (630 links) achieve 100\% survival, but require $1.8\times$ and $2.7\times$ as much cabling cost and demand high port counts on core switches. Our hybrid dual-plane layout is the cheapest graph (102 links, cost 234) that provides two node-disjoint paths and survives 100\% of single and double faults.

\begin{table}[t]
\caption{Topology comparison (same campus; selected architectures).}
\label{tab:topo}
\centering
\footnotesize
\begin{tabular}{@{}lrrrrr@{}}
\toprule
Topology & Links & Cost & Hops & Disj.\ paths & 2-fail surv. \\
\midrule
Hybrid dual-plane (ours) & 102 & 234 & 3.71 & 2.0 & 100\% \\
3-tier single-homed      & 49  & 111 & 3.71 & 1.0 & 92\% \\
Distribution ring        & 42  & 84  & 3.54 & 1.0 & 90\% \\
Spine-leaf (Clos)        & 144 & 432 & 2.00 & 4.0 & 100\% \\
Full mesh (reference)    & 630 & 630 & 1.00 & 35  & 100\% \\
\bottomrule
\end{tabular}
\end{table}

%==============================================================================
\section{Discussion, Limitations, and Future Work}\label{sec:discuss}
%==============================================================================
Building and testing the fabric brought out three practical takeaways. First, many SDN papers claim sub-50\,ms recovery on Mininet without verifying whether the software switch actually flipped its local FF bucket. Checking group packet counters showed that software OVS leaves the backup bucket idle on a virtual link-down unless BFD is active, meaning recovery is actually driven by the controller (80--210\,ms). Reporting that behavior plainly makes our numbers reproducible. Second, our control primitive works only because the physical wiring backs it up. First-hop plane selection is safe because Plane~A and Plane~B share zero transit switches; on a random mesh, flipping an ingress bucket could easily dump backup packets into a downstream failed switch. Third, handling protection, WCMP load spreading, and DiffServ EF pinning inside the same first-hop group entry keeps switch table usage tiny ($\approx$43 rules per access switch) instead of layering STP, OSPF, and MPLS on top of each other.

Several scope boundaries remain for future work. Because our timings come from a software Mininet/OVS testbed on one machine, absolute hardware timings will differ; moving the controller onto physical OpenFlow switches with hardware BFD is the next step to measure true sub-50\,ms local FF cutover. Forwarding currently matches Layer~2 MAC prefixes, so adding ARP handling and Layer~3 subnet routing is needed for multi-VLAN campuses. We also verified WCMP weights and HTB queues functionally, but we still need to run sustained \texttt{iperf3} congestion benchmarks to log tail-latency curves under heavy load. Finally, while our two-node REST heartbeat script works well for a single campus, multi-campus setups would benefit from a consensus-backed store.

%==============================================================================
\section{Conclusion}\label{sec:conclusion}
%==============================================================================
We built and verified a campus SDN fabric where fast failover, multipath load balancing, and DiffServ QoS reduce to a single first-hop plane-selection choice over two node-disjoint planes on OpenFlow~1.3. Proactive destination-rooted group tables keep all physical links active and recovered 100\% of tested link and switch faults in 80--210\,ms with minimal controller CPU and RAM load. An offline prover across all 510 single-link and 25{,}755 double-link fault cases confirmed zero loops or black holes whenever a physical path exists, and our graph comparison showed that the dual-plane wiring is the lowest-cost layout to survive every single and double fault. Paired with a primary-backup controller heartbeat daemon, the design replaces three separate legacy protocols with one verifiable OpenFlow primitive.

%==============================================================================
\begin{thebibliography}{00}
\bibitem{sharma} S.~Sharma, D.~Staessens, D.~Colle, M.~Pickavet, and P.~Demeester,
``OpenFlow: Meeting carrier-grade recovery requirements,'' \emph{Computer
Communications}, vol.~36, no.~6, pp.~656--665, 2013.
\bibitem{isyaku} B.~Isyaku \emph{et al.}, ``Dynamic Routing and Failure Recovery
Approaches for Efficient Resource Utilization in OpenFlow-SDN: A Survey,''
\emph{IEEE Access}, vol.~10, 2022.
\bibitem{menaceur} ``Fault Tolerance and Failure Recovery Techniques in
Software-Defined Networking: A Comprehensive Approach,'' \emph{Journal of
Network and Systems Management}, 2023.
\bibitem{fapr} ``FAPR: An Adaptive Approach to Link Failure Recovery in SDN with
High Speed and Low Interruption Rate,'' \emph{Applied Sciences (MDPI)}, 2024.
\bibitem{ijaip} ``OpenFlow groups based fast failover mechanism for
software-defined networks,'' \emph{Int.\ J.\ of Advanced Intelligence
Paradigms}, vol.~28, 2024.
\bibitem{survey_lf} ``A Survey of Link Failure Detection and Recovery in
Software-Defined Networks,'' \emph{Computer Systems Science and Engineering},
2024.
\bibitem{safelb} ``Safe Load Balancing in Software-Defined Networking,'' in
\emph{Proc.\ IEEE/IFIP NOMS}, 2024.
\bibitem{transformer} ``A Transformer-based Deep Q-Learning Approach for Dynamic
Load Balancing in Software-Defined Networks,'' preprint, 2024.
\bibitem{ha2024} ``Fault-tolerant distributed SDN control plane,'' \emph{Computer
Networks (Elsevier)}, vol.~241, 2024.
\bibitem{openflow} Open Networking Foundation, ``OpenFlow Switch Specification,
Version 1.3,'' 2012.
\bibitem{mininet} B.~Lantz, B.~Heller, and N.~McKeown, ``A network in a laptop:
Rapid prototyping for software-defined networks,'' in \emph{Proc.\ ACM HotNets},
2010.
\bibitem{ovs} B.~Pfaff \emph{et al.}, ``The design and implementation of Open
vSwitch,'' in \emph{Proc.\ USENIX NSDI}, 2015.
\bibitem{ryu} Ryu SDN Framework Community, ``Ryu: Component-based software-defined
networking framework,'' \url{https://ryu-sdn.org}.
\bibitem{wcmp} J.~Zhou \emph{et al.}, ``WCMP: Weighted cost multipathing for
improved fairness in data centers,'' in \emph{Proc.\ ACM EuroSys}, 2014.
\bibitem{diffserv} K.~Nichols, S.~Blake, F.~Baker, and D.~Black, ``Definition of
the Differentiated Services Field (DS Field) in the IPv4 and IPv6 Headers,'' RFC
2474, 1998.
\bibitem{bfd} D.~Katz and D.~Ward, ``Bidirectional Forwarding Detection (BFD),''
RFC 5880, 2010.
\end{thebibliography}

\end{document}
